\chapter{Retrieval and Grounding}
\label{chapter:Retrieval}
% EVOLVE-BLOCK-START

An agent knows only what it reads.\index{Retrieval}
A language model's parametric knowledge, the information encoded in its weights during training, is broad but limited.
It can be generic, stale, and unreliable with respect to the specific details that matter for any particular research task.
A physicist's agent should know about dimensional analysis; a biologist's should know about enzyme kinetics; an economist's should know about functional forms in demand estimation.
None of this domain-specific knowledge can be assumed to be reliably present in the model's weights.

\emph{Retrieval-augmented generation} (RAG)\index{Retrieval-augmented generation} (Lewis et al., 2020) addresses this by injecting relevant domain knowledge into the context at the time it is needed.
Rather than relying on the model to remember a fact, the loop retrieves the fact from a curated corpus and places it directly in the window.

In the course, retrieval is not a secondary feature; it has a real job.
The reference corpus holds domain units, functional forms, and dimensional constraints that constrain which model proposals are physically plausible.
The challenge is to demonstrate that retrieval moves the leaderboard: a grounded agent should propose better candidates because its proposals are constrained by domain knowledge.

\section{The Retrieval Problem}
\label{section:RetrievalProblem}

\subsection{The Knowledge Gap}

Let $\mathcal{K}$ denote the body of domain knowledge relevant to a task (scientific literature, textbooks, experimental notes).
The context window can hold only a small fraction of $\mathcal{K}$ at any invocation.
The \emph{retrieval problem} is to select the subset of $\mathcal{K}$ most relevant to the current query and place it in the context.
Formally, given a query $q$ (the agent's current task description or question) and a corpus $\mathcal{K}$, return the $k$ passages from $\mathcal{K}$ that are most relevant to $q$.
In practice $k$ is small, a handful of passages rather than hundreds, because retrieved text competes for the same window budget as the agent's instructions, its tools, and its accumulated history.

\begin{remark}
Retrieval is not automatically the right tool.
When the entire relevant corpus fits comfortably in the context window, a threshold on the order of a few hundred pages for current models, loading it directly can beat retrieval on questions that require synthesis across the whole document (Chapter~\ref{chapter:ContextEngineering}), because no relevant passage can be missed.
Retrieval earns its place when $\mathcal{K}$ is far larger than any window, when only a small slice is relevant per query, or when injecting irrelevant material would crowd out what matters.
The course corpora are deliberately small, so this boundary is worth checking before building an index at all.
\end{remark}

\subsection{Keyword vs.\ Semantic Retrieval}

Classical retrieval systems (e.g., BM25, TF-IDF) match keywords: they return passages that contain the same words as the query.\index{Keyword retrieval}
These systems fail when the query and the relevant passage use different vocabulary, for example, when the query asks about ``drag force'' and the relevant passage discusses ``aerodynamic resistance.''

\emph{Semantic retrieval}\index{Semantic retrieval} avoids this failure by representing both query and passage as vectors in a shared semantic space, and returning the passages whose vectors are closest to the query vector.
Two passages about the same concept have similar vectors even if they use different words.
The word \emph{shared} carries the weight here: the query and every passage are mapped by the same function into the same space, so their coordinates are directly comparable and geometric closeness can stand in for closeness in meaning.
In current practice the two are usually combined rather than chosen between: a \emph{hybrid} index runs keyword and semantic retrieval together and merges their results, and the contextual-retrieval technique of Section~\ref{section:ContextualRetrieval} sharpens both at once.
Keyword matching keeps its value precisely where meaning blurs: an exact identifier such as a gene name, a statute number, or the symbol one particular paper uses for a parameter is matched more reliably by the literal string than by any embedding of it.

\section{Embeddings and Semantic Similarity}
\label{section:Embeddings}

\subsection{What an Embedding Is}

An \emph{embedding}\index{Embedding} is a function $\phi: V^* \to \RealNumbers^d$ that maps a text passage to a real vector.
Embedding models are trained so that semantically similar passages map to nearby vectors under a standard distance metric.

Two pieces of that signature deserve unpacking.
The domain $V^*$ is the set of all finite strings of text, so $\phi$ accepts a five-word query and a five-hundred-word passage alike, while the codomain $\RealNumbers^d$ has a \emph{fixed} dimension, typically between~768 and~3072 in current models, so both inputs come back as vectors of the same length and can therefore be compared.
The function itself is a neural network trained on very large collections of text pairs known to belong together (a question and the passage that answers it, a title and its abstract), under an objective that pulls the vectors of related pairs together and pushes unrelated pairs apart.
Its coordinates are not individually interpretable: no axis means ``biology,'' no single number can be read on its own, and only the relative geometry of the vectors carries meaning.
In practice a researcher never touches the network: producing an embedding is one API call that takes a string and returns a list of $d$ floating-point numbers.

\begin{definition}
\label{definition:Embedding}
An \emph{embedding model} is a function $\phi: V^* \to \RealNumbers^d$ trained so that, for \emph{most} triples of passages $p_1, p_2, p_3$ in which $p_1$ is more semantically similar to $p_2$ than to $p_3$,
\begin{equation}
\mathrm{sim}(\phi(p_1), \phi(p_2)) > \mathrm{sim}(\phi(p_1), \phi(p_3)) ,
\end{equation}
where $\mathrm{sim}$ is a similarity function, typically cosine similarity.
The property is statistical rather than universal: no finite-dimensional embedding orders \emph{every} triple correctly, and retrieval quality depends on how often the inequality holds for the passages that matter.
\end{definition}

The definition is deliberately a statement about \emph{ordering} rather than about distances: it constrains which of two passages ranks closer to a third, and says nothing about what any particular similarity value means in absolute terms.
That is exactly the guarantee retrieval needs, because a retrieval system only ever sorts candidates and keeps the top few.

\subsection{Cosine Similarity}

\emph{Cosine similarity}\index{Cosine similarity} measures the angle between two vectors, independent of their magnitude:
\begin{equation}
\mathrm{sim}(u, v) = \frac{u \cdot v}{\|u\| \, \|v\|} \in [-1, 1] .
\label{equation:CosineSimilarity}
\end{equation}
Values near~1 indicate high similarity; values near~0 indicate orthogonality (unrelated topics); values near $-1$ occur rarely and indicate strong opposition.
In practice the cosine similarities produced by a trained text-embedding model occupy a compressed positive band rather than the full $[-1, 1]$ range: unrelated passages typically score well above~0, so a value like~$0.23$ should be read as ``much less similar,'' not as literal orthogonality.
Cosine similarity is preferred over Euclidean distance because it is invariant to the magnitude of the embedding vectors, which vary with passage length and model implementation.
The practical consequence of the compressed band is that absolute cutoffs do not transfer: a rule such as ``keep everything scoring above~$0.8$'' tuned on one embedding model or one corpus is meaningless on another, so either calibrate the cutoff against the score distribution of your own corpus or avoid cutoffs entirely and simply take the top~$k$.

\begin{example}
Embedding two domain passages and computing their similarity:
\begin{lstlisting}[language=python]
import numpy as np

def cosine_similarity(a, b):
    return np.dot(a, b) / (np.linalg.norm(a) * np.linalg.norm(b))

# embed(): stand-in for an embedding API call (cf. embed_fn below);
# the printed values are illustrative, not computed here (d = 1536)
emb_drag    = embed("Drag force scales as the square of velocity")
emb_resist  = embed("Aerodynamic resistance increases with v^2")
emb_density = embed("Material density is mass divided by volume")

print(cosine_similarity(emb_drag, emb_resist))   # ~0.91 (illustrative) — similar
print(cosine_similarity(emb_drag, emb_density))  # ~0.23 (illustrative) — much less similar
\end{lstlisting}
What to notice is the pair that scores high: those two passages share no content word at all (``drag force'' against ``aerodynamic resistance,'' ``the square of velocity'' against ``$v^2$''), so a keyword index would rank them as unrelated, while the embedding places them beside each other.
\end{example}

\section{Vector Search}
\label{section:VectorSearch}

\subsection{The Nearest-Neighbor Problem}

Given a corpus of $N$ embedded passages $\{\phi(p_i)\}_{i=1}^N$ and a query embedding $\phi(q)$, the retrieval problem reduces to finding the $k$ passages with the highest cosine similarity to $\phi(q)$.
This is the \emph{$k$-nearest-neighbor problem} in $\RealNumbers^d$.\index{$k$-nearest-neighbor}

Exact nearest-neighbor search is $O(N \cdot d)$ per query, linear in the corpus size.
Concretely, one similarity score costs $d$ multiplications and additions, so a corpus of ten thousand passages at $d = 1536$ costs roughly fifteen million multiply-adds per query: milliseconds on a laptop.
For the small corpora used in the course (hundreds to thousands of passages), exact search is fast enough.
For larger corpora, \emph{approximate nearest-neighbor} (ANN) indices trade a small reduction in recall for a large reduction in query time (Johnson et al., 2021).
Recall here is the fraction of the true top~$k$ that the approximate search actually returns, and the speed is bought by grouping vectors into neighborhoods and searching only the promising ones, accepting that the search will occasionally walk past a genuine neighbor.

\subsection{Building a Retrieval Index}

The workflow for building a retrieval system has three stages:
\begin{enumerate}
\item \textbf{Chunk.} Split the corpus into passages of a manageable size (typically 200--500 tokens).
Those counts amount to roughly one to three paragraphs of English prose.
Smaller chunks retrieve more precisely; larger chunks provide more context per retrieved passage.
The choice matters more than its mechanical description suggests, because the chunk is the unit of retrieval: whatever the index returns, it returns whole chunks, so chunk size fixes the granularity of what the agent can be handed.
It is equally the unit of representation, since a single vector must summarize the entire chunk; a forty-page paper embedded whole yields one point averaging its introduction, its methods, and its discussion, sitting vaguely near everything and decisively near nothing.

\item \textbf{Embed.} Run each passage through the embedding model to produce a vector.
Store the vectors alongside the passage text.

\item \textbf{Index.} Build a data structure (a flat array for exact search, or an ANN index for approximate search) that supports efficient nearest-neighbor queries.
The word \emph{index} promises something more exotic than it is: for exact search it is literally the $N \times d$ array of stored vectors plus the scan that reads it, and it earns the name only once the structure starts skipping candidates.
\end{enumerate}

At query time: embed the query, search the index for the $k$ nearest passage vectors, retrieve the corresponding texts, and inject them into the agent's context.

\begin{example}
A minimal retrieval index using NumPy:
\begin{lstlisting}[language=python]
class FlatIndex:
    # embed_fn: any embedding API call, e.g.
    #   lambda p: client.embeddings.create(input=p, model=...).data[0].embedding
    def __init__(self, passages, embed_fn):
        self.passages = passages
        self.vectors  = np.array([embed_fn(p) for p in passages])
        # normalise for cosine similarity via inner product
        norms = np.linalg.norm(self.vectors, axis=1, keepdims=True)
        norms = np.maximum(norms, 1e-12)   # guard against zero-norm embeddings
        self.vectors = self.vectors / norms

    def query(self, q_vec, k=5):
        q_norm = q_vec / np.maximum(np.linalg.norm(q_vec), 1e-12)
        scores = self.vectors @ q_norm
        top_k  = np.argsort(scores)[::-1][:k]
        return [(self.passages[i], float(scores[i])) for i in top_k]
\end{lstlisting}
Two details in this listing carry the idea.
Every stored vector is normalized once at build time, which turns Eq.~\eqref{equation:CosineSimilarity} into a plain inner product, so the whole search collapses to the single line \texttt{self.vectors @ q\_norm}: one matrix-vector product against the entire corpus.
And the passage text is stored alongside the vectors, because what the agent eventually reads is text; the vectors exist only to decide \emph{which} text.
\end{example}

\subsection{Contextual Retrieval}
\label{section:ContextualRetrieval}

A chunk extracted in isolation loses the context of the document it came from.
A passage reading ``the growth rate falls to zero at high density'' does not say which population, which model, or under what assumption; and the embedding of that fragment is correspondingly ambiguous.
\emph{Contextual retrieval}\index{Contextual retrieval} (Anthropic, 2024) repairs this by using a language model to prepend a short, chunk-specific description of the surrounding document to each chunk \emph{before} it is embedded and indexed.
Concretely, the ambiguous fragment above might be stored and embedded as ``From a chapter on logistic population growth, describing the model's behavior as the population $P$ approaches the carrying capacity $K$: the growth rate falls to zero at high density.''
The added clause costs a few dozen tokens and moves the chunk's vector next to the queries about carrying capacity that the bare fragment would have missed.
The enriched chunk feeds both the semantic index and a keyword index, so semantic and lexical matching reinforce each other rather than compete.
On the source's benchmark, contextualizing the embeddings reduced retrieval failures by roughly a third, and adding contextual keyword matching by roughly a half; the technique modifies the \textbf{embed} stage of the pipeline above, not the architecture around it.
These reductions compose, as Figure~\ref{figure:RetrievalStaircase} shows, which is the point: retrieval quality is a measurable quantity that improves in stackable increments, not a single knob to be tuned by feel.

\begin{figure}[htb!]
\begin{center}
%% retrieval-staircase.tex: retrieval failure rate falling in stackable increments
%% as techniques are layered (Chapter 8).
%% A data figure, not a block diagram: hand-drawn bars against a baseline rule, so
%% the house conceptual styles do not apply.
\begin{tikzpicture}[arr/.style={->, >=stealth', thick}]
\draw[arr] (-0.3,0) -- (-0.3,4.1) node[above, font=\scriptsize] {failure rate};
\draw[arr] (-0.3,0) -- (11.4,0) node[right, font=\scriptsize] {stages added};
\draw[dashed] (-0.3,3.6) -- (10.9,3.6) node[right, font=\scriptsize] {baseline};
\draw[fill=gray!15]   (0.5,0)  rectangle (2.3,3.6);
\draw[fill=blue!15]   (3.3,0)  rectangle (5.1,2.4);
\draw[fill=green!20]  (6.1,0)  rectangle (7.9,1.8);
\draw[fill=yellow!20] (8.9,0)  rectangle (10.7,1.2);
\node[font=\scriptsize] at (4.2,2.75) {$\downarrow \tfrac{1}{3}$};
\node[font=\scriptsize] at (7.0,2.15) {$\downarrow \tfrac{1}{2}$};
\node[font=\scriptsize] at (9.8,1.55) {$\downarrow \tfrac{2}{3}$};
\node[font=\scriptsize, align=center] at (1.4,-0.5) {baseline};
\node[font=\scriptsize, align=center] at (4.2,-0.5) {$+$ contextual\\embeddings};
\node[font=\scriptsize, align=center] at (7.0,-0.5) {$+$ contextual\\keyword};
\node[font=\scriptsize, align=center] at (9.8,-0.5) {$+$ rerank};
\end{tikzpicture}%

\caption{Retrieval failure rate (the fraction of queries whose relevant passage is missing from the top~$k$) falling in stackable increments as techniques are layered, on the source's benchmark (Anthropic, 2024): contextual embeddings alone cut baseline failures by roughly a third, adding contextual keyword matching reaches roughly a half, and a reranking second stage reaches roughly two-thirds.
The takeaway is not the exact numbers, which are benchmark-specific and will shift, but the shape: each stage lowers the same measurable quantity, so retrieval quality is evaluated and improved the way Chapter~\ref{chapter:Evaluation} treats any score, not judged by feel.}
\label{figure:RetrievalStaircase}
\end{center}
\end{figure}

\begin{remark}
Retrieval quality improves further with a second stage.
A \emph{reranker}\index{Reranker} takes a generous first-pass candidate set, say the top~150 by vector similarity, and rescores each candidate against the query with a more expensive model, keeping only the best handful.
That second model is more accurate because it reads the query and the candidate passage \emph{together} and scores the pair, instead of embedding each separately and comparing afterwards.
Such a joint reading cannot be precomputed, which is precisely why it is unaffordable across a whole corpus and is confined to a shortlist.
The pattern is general: a cheap, high-recall first stage followed by a precise, low-throughput second stage, which on the source's benchmark cut retrieval failures by about two-thirds overall (Anthropic, 2024).
\end{remark}

\begin{remark}
Generating a context sentence for every chunk sounds expensive, but \emph{prompt caching} makes it cheap: the parent document is cached once and reused across all of its chunks rather than re-sent per chunk.
Prompt caching is a provider-side facility that retains the processed form of a repeated prompt prefix, so each subsequent call beginning with the same document pays a reduced rate for that prefix instead of the full input price.
The source reports a one-time cost on the order of a dollar per million document tokens (Anthropic, 2024), small enough that contextualization is a routine preprocessing step rather than a luxury.
\end{remark}

\begin{remark}
Flat similarity search is not the only shape a retrieval index can take.
When the relationships \emph{between} entities carry more of the answer than the similarity of any single passage (which protein interacts with which, which result cites which) a \emph{structural} approach can help: build a graph over entities and their relations extracted from the corpus and retrieve connected subgraphs rather than isolated chunks, the family of techniques often called \emph{GraphRAG}.
The tradeoff is a heavier, more brittle indexing step (entity and relation extraction can err) against better recall on multi-hop questions that a flat index answers poorly.
As with grounding itself, the honest test is the ablation of Section~\ref{section:GroundingConstraint}: adopt the more complex index only where it measurably beats flat vector search on the questions that matter.
\end{remark}

\section{Retrieval-Augmented Research}
\label{section:RAG}

\subsection{The RAG Architecture}

A \emph{retrieval-augmented} agent extends the basic loop with a retrieval step before each model invocation (Figure~\ref{figure:RAGPipeline}).\index{RAG architecture}

\begin{enumerate}
\item The agent (or the loop) forms a query from the current task: ``What functional forms are commonly used to model population growth?''
\item The query is embedded and passed to the retrieval index.
\item The $k$ most relevant passages are retrieved from the corpus.
\item The passages are injected into the agent's context as domain knowledge.
\item The agent reads the augmented context and produces a proposal conditioned on the retrieved information.
\end{enumerate}

\begin{figure}[htb!]
\begin{center}
%% rag-pipeline.tex: the retrieval-augmented generation pipeline (Chapter 8).
%% House conceptual style; the offline, sealed corpus is drawn as a dashed `store'.
%% _concept-style.tex: shared TikZ styles for the course's CONCEPTUAL block diagrams.
%% The leading underscore marks this as the one file in figures/ that is not a
%% figure; every other figures/*.tex holds exactly one figure body.
%%
%% \input this at the top of every conceptual figure (before \begin{tikzpicture})
%% so each such diagram speaks one visual language: rounded "block" nodes, stealth
%% arrows, and natural `<hue>!15' fills. Redefining these styles on each \input is
%% idempotent and harmless. The reference for the house parameters is
%% resources/STYLE-figures.md. Figure~\ref{figure:AgentLoop} is the exemplar.
%%
%% Scope: use for conceptual/relational diagrams (boxes-and-arrows). Do NOT use for
%% product mock-ups (the rig figures have their own shared language) or for
%% data/plot figures.
\tikzset{
  % the standard block node: geometry only — apply a `fill=<hue>!15' per node.
  conceptbox/.style={draw, rounded corners=4pt, minimum width=2.4cm,
                     minimum height=0.85cm, align=center, font=\small},
  % a flow arrow between blocks.
  conceptflow/.style={->, >=stealth', thick},
  % an edge/annotation label.
  conceptlabel/.style={font=\scriptsize, align=center},
}%

\begin{tikzpicture}[
    node distance=1cm,
    box/.style={conceptbox, minimum width=2.2cm, minimum height=1.0cm},
    store/.style={conceptbox, dashed, fill=gray!15,
                  minimum width=3.0cm, minimum height=1.0cm},
    arr/.style={conceptflow}
]
\node[box, fill=blue!15]                   (q)     {query $q$};
\node[box, fill=orange!15, below=of q]      (qe)    {embed\\$\phi(q)$};
\node[box, fill=green!15,  right=of qe]     (idx)   {vector index\\search};
\node[box, fill=green!15,  right=of idx]    (topk)  {top-$k$\\passages};
\node[box, fill=violet!15, right=of topk]   (model) {model reads\\augmented context};
\node[store, above=of idx]            (corpus) {corpus $\mathcal{K}$ (sealed)\\chunk $+$ embed, offline};

\draw[arr] (q)    -- (qe);
\draw[arr] (qe)   -- node[above,font=\scriptsize]{$\phi(q)$} (idx);
\draw[arr] (idx)  -- node[above,font=\scriptsize]{$k$-NN}    (topk);
\draw[arr] (topk) -- node[above,font=\scriptsize]{inject}    (model);
\draw[arr] (corpus) -- (idx);
\end{tikzpicture}%

\caption{The retrieval-augmented generation pipeline.
At query time, the query is embedded, matched against the vector index by $k$-nearest-neighbor search, and its top-$k$ passages are injected into the context the model reads.
The corpus $\mathcal{K}$ is chunked and embedded once, offline (dashed), and sealed thereafter; exactly the discipline of Section~\ref{section:VectorSearch}.
Retrieval quality is set at the index-search step: an index that returns irrelevant passages spends context budget for no gain, and one that returns misleading passages actively steers the proposal wrong.}
\label{figure:RAGPipeline}
\end{center}
\end{figure}

The quality of the retrieval determines the quality of the grounding.
A retrieval system that returns irrelevant passages wastes context budget; one that returns the wrong information can actively mislead the agent.

\subsection{What Goes in the Corpus}

The reference corpus used for model discovery in this course contains domain notes organized by physical or scientific domain: common functional forms, dimensional analysis rules, typical parameter ranges, and known constraints.

\begin{example}
A representative corpus entry for population dynamics:
\begin{lstlisting}[language=markdown]
## Logistic Growth

Population dynamics frequently follow logistic growth when a carrying
capacity limits the population. The standard form is:

    P(t) = K / (1 + ((K - P0) / P0) * exp(-r * t))

where K is the carrying capacity, P0 is the initial population,
r > 0 is the intrinsic growth rate, and t is time.

At low densities (P << K), logistic growth approximates exponential
growth. At high densities, the growth rate decreases to zero.

**Dimensional constraint**: r has units of 1/time; K and P have the
same units (individuals, biomass, etc.).
\end{lstlisting}
When the agent is given data from a population dynamics experiment, retrieving this passage provides exactly the prior knowledge needed to constrain proposals to plausible functional forms.
Notice that the entry is written to be found: a heading that names the concept, the closed-form expression, the limiting behavior at both extremes, and an explicit dimensional-constraint line.
Each of those is a hook a query can match, and each is plain text rather than a figure or a typeset table, because the chunk that gets retrieved is exactly the text the model will read.
\end{example}

\section{Grounding as Constraint}
\label{section:GroundingConstraint}

The central claim of this chapter is that domain grounding \emph{moves the leaderboard}: a retrieval-augmented agent proposes better candidates because it is constrained to the space of scientifically plausible models.

\subsection{Constraint Mechanisms}

Grounding constrains model proposals in two ways.

\begin{enumerate}
\item \textbf{Form seeding.} Retrieved functional forms (logistic, power-law, exponential decay) seed the proposal distribution.
Instead of proposing arbitrary polynomial combinations, the agent proposes variations of known-plausible forms.

\item \textbf{Dimensional filtering.} Retrieved dimensional constraints rule out proposals that are dimensionally inconsistent.
If the target is a rate with units $1/\text{time}$, and a proposed model has units $\text{mass}/\text{time}^2$, the mismatch can be detected and the proposal discarded before any fitting is attempted.
\end{enumerate}

\begin{example}
The constraint line of the logistic corpus entry above supports a filter a reader can apply by hand to a proposed rate law.
That line states that $r$ has units of $1/\text{time}$ and that $K$ carries the same units as the population $P$.
The proposal $\dot{P} = r P (1 - P/K)$ passes, since $P/K$ is dimensionless and both sides are therefore a population per unit time.
The proposal $\dot{P} = r P (1 - P) - K$ fails twice over: $1 - P$ subtracts a population from a pure number, and the trailing $K$ adds a population to a rate.
\end{example}

Together, these mechanisms shrink the search space to a plausible region.
A search over a smaller, better-chosen space outperforms a search over a larger, undifferentiated space (provided the constraints are correct).

\subsection{When Grounding Helps and When It Hurts}

Grounding is not always beneficial.
It helps when the retrieved knowledge is relevant and the domain is well-understood; it hurts when the retrieved knowledge is irrelevant, misleading, or too strong a prior.

\begin{example}
A dataset generated from the formula $y = e^{-x_1^2} \sin(x_2)$ is unusual (a product of Gaussian decay and oscillation, a functional form that rarely appears in domain corpora organized around named physical laws).
If the corpus contains no entries about oscillatory-times-Gaussian forms, retrieval will return the nearest available match (perhaps a simple exponential decay), which could steer the agent away from the correct form.
In this case, the ungrounded agent performs better than the grounded one because the domain prior is misaligned.
The mechanism is worth naming: a similarity index has no way to report that nothing it holds is relevant, since it ranks whatever it has and returns the top~$k$ regardless, so an absent answer arrives looking exactly like a present one.
Guarding against that takes an explicit score threshold or a reranker permitted to hand back an empty shortlist.
\end{example}

\begin{remark}
Grounding helps when the retrieved prior is aligned and hurts when it is misaligned, and which case obtains cannot be assumed in advance.
The grounded versus ungrounded ablation required in this course, run on at least two datasets, exists precisely to quantify this tradeoff rather than presume it.
A dataset where grounding helps and a dataset where it hurts, analyzed together, build the researcher's intuition for when retrieval adds value.
\end{remark}

\section*{Further Reading}

\emph{The engineering-blog and product references below are Fall~2026 snapshots of a fast-moving surface; titles and locations may have since changed.}

\begin{small}
\begin{enumerate}
\item Lewis, P., et al., ``Retrieval-Augmented Generation for Knowledge-Intensive NLP Tasks,'' \emph{NeurIPS}, 2020 (arXiv:2005.11401) — the original RAG paper.
\item Gao, Y., et al., ``Retrieval-Augmented Generation for Large Language Models: A Survey,'' arXiv:2312.10997, 2023 — a comprehensive survey of RAG architectures and applications.
\item Johnson, J., Douze, M., and Jégou, H., ``Billion-Scale Similarity Search with GPUs,'' \emph{IEEE Transactions on Big Data}, 2021 (arXiv:1702.08734) — the paper behind FAISS, the most widely used ANN library.
\item Anthropic, ``Introducing Contextual Retrieval,'' 2024 — prepending chunk-specific context before embedding and keyword indexing, with reranking; the source for Section~\ref{section:ContextualRetrieval}.
\item Edge, D., et al., ``From Local to Global: A Graph RAG Approach to Query-Focused Summarization,'' \emph{Microsoft Research}, 2024 — the structural-retrieval alternative: build a graph over extracted entities and relations and retrieve connected subgraphs, useful when relationships carry the answer.
\end{enumerate}
\end{small}
% EVOLVE-BLOCK-END
